\section*{Ethical Considerations}
\label{sec:ethics}
In this section, we discuss the ethical considerations of our study and how they were addressed in our implementation of \citeb and the associated measurements.

\noindent\textbf{Responsible Use of External Services.}
\label{sec:ethical_use_external_services}
All LLMs and external services were accessed in compliance with their respective terms of service and usage guidelines.
We applied rate limiting and conservative request schedules to avoid overloading provider infrastructure, and we did not use adversarial or harmful prompts, following the Menlo report~\cite{bailey2012menlo}.

\noindent\textbf{Data Access and Stewardship.}
\label{sec:data_stewardship}
Paper collection relied on publicly accessible proceedings and institutionally authorized access by our research team.
For open-access venues, we used DBLP to obtain DOI links and downloaded papers via automated scripts; for venues that provide full proceedings, we downloaded the proceedings and split PDFs by the table of contents; for access-restricted venues, we used automated browser control to mimic manual downloads under institutional access.
We enforced a 5-second delay between downloads to avoid server burden, and any access-restricted papers were retrieved only through the researchers' institutional access rights.
All papers are used solely for academic research, and no portion of the corpus has been redistributed or disseminated to the outside.

\noindent\textbf{Human Subjects and Privacy.}
The survey received IRB approval from our institution prior to data collection.
Participants were recruited through two transparent channels: we published the survey link on social media platforms, and we randomly sampled 300 researchers from program committees and author lists of top-tier venues for direct email outreach through their institutional or publicly available email addresses.
The survey provided full transparency to participants, including detailed explanation of the research purpose, methodology, contact information for the research team, and an explicit opt-out mechanism allowing withdrawal at any point.
The questionnaire did not solicit personally identifiable information; all responses were stored securely in de-identified form and have not been shared outside the research.
% team.

\noindent\textbf{Anonymization of Papers with Invalid Citations.}
Following the precedent of meta-science studies that critically examine published work~\cite{schloegel2025confusing}, we decided not to disclose the identities of specific papers containing invalid citations in the submitted manuscript.
Our goal is not to blame individual authors, institutions, or papers, but to inform and spark constructive discussion around a widespread phenomenon: the presence of ghost citations (untraceable or fabricated references) in the scientific record, that affects the entire research community.
Invalid citations may arise for various reasons (e.g., researcher oversight, uncritical adoption of LLM-generated content, or gaps in verification workflows) and do not necessarily constitute evidence of intent.
We acknowledge that causality and attribution are complex; our focus is on systemic patterns rather than individual cases.
We welcome feedback from reviewers on how best to balance transparency with proportionality in addressing this emerging challenge.
